% =========================================================================
% English translation (generated by AI using Claude Opus 4.8 (Max effort)).
% This file was translated from Hungarian to English by AI.
% DISCLAIMER: Please review against the original Hungarian .tex files, before relying on it!
% SOURCE: 1. Függvények határértéke, folytonossága.tex
% =========================================================================
\documentclass[margin=0px]{article}

\usepackage{listings}
\usepackage[utf8]{inputenc}
\usepackage{graphicx}
\usepackage{float}
\usepackage[a4paper, margin=0.7in]{geometry}
\usepackage{amsthm}
\usepackage{amssymb}
\usepackage{amsmath}
\usepackage{fancyhdr}
\usepackage{setspace}

\onehalfspacing

\newenvironment{tetel}[1]{\paragraph{#1 \\}}{}

\pagestyle{fancy}
\lhead{\it{CS BSc Final Exam Topics}}
\rhead{1. Limits and continuity of functions}

\title{\textbf{{\Large ELTE IK - Computer Science BSc} \vspace{0.2cm} \\ {\huge Final Exam Topics}} \vspace{0.3cm} \\ 1. Limits and continuity of functions}
\author{}
\date{}

\begin{document}
\maketitle

\begin{center}
\fbox{\parbox{0.92\textwidth}{\small \textit{Note: This document was translated from Hungarian to English by AI. Please verify the information against the original Hungarian source before relying on it.}}}
\end{center}

\begin{tetel}{Limits and continuity of functions}
    Limits of functions, continuity. Properties of functions that are continuous on a compact set: Weierstrass theorem, Bolzano theorem. The concept of a power series, Cauchy--Hadamard theorem, analytic functions.
\end{tetel}

\section{Limits of functions}

Given $f \in \mathbb{K}_{1} \to \mathbb{K}_{2}, a \in \mathcal{D}'_{f}$ (accumulation point). The function $f$ has a limit at the point $a$ if $\exists A \in \overline{\mathbb{K}}_{2}$, where $\overline{\mathbb{K}} = \mathbb{C} \vee \mathbb{R} \vee +\infty \vee -\infty$, such that for any neighborhood $K(A) \subset \mathbb{K}_{2}$ a suitable neighborhood $k(a) \subset \mathbb{K}_{1}$ of $a$ can be given with which $f[(k(a) \backslash \{a\}) \cap \mathcal{D}_{f}] \subset K(A)$ holds. \\
Otherwise: $f(x) \in K(A), (a \neq x \in k(a) \cap \mathcal{D}_{f})$. \\
Then the uniquely existing number $A \in \overline{\mathbb{K}}_{2}$ or one of the infinities is called the limit of the function $f$ at the point $a$. \\
Notation: $\lim\limits_{x \to a}{f(x)} = \lim\limits_{a}{f} = A$

\subsection{Accumulation point}

$A \subset \mathbb{K}$, then $\alpha \in \overline{\mathbb{K}}$ is an accumulation point of the set $A$ if for any neighborhood $K(\alpha)$ we have $(K(\alpha) \backslash \{\alpha\}) \cap A \neq \emptyset$. \\
With inequalities: $A \subset \mathbb{K}$, then the number $\alpha \in \mathbb{K}$ is an accumulation point of the set $A$ if for $\forall \varepsilon > 0$ there exists $x \in A$ such that $0 < |x-\alpha| < \varepsilon$.

\subsection{Neighborhood}

$A \subset \mathbb{K}, a \in A, r > 0: K_{r}(a) = K(a) = \{x \in A: |x-a| < r\}$.

\subsection{Function}

$\emptyset \neq A, B$ sets, $f \subset A \times B$ a relation. For some element $x \in A$ let $f_{x} := \{y \in B: (x,y) \in f\}$. The relation $f$ is a function if for $\forall x \in A$ the set $f_{x}$ has at most one element.

\section{Continuity of functions}

The function $f \in \mathbb{K}_{1} \to \mathbb{K}_{2}$ is continuous at the point $a \in \mathcal{D}_{f}$ if for $\forall \varepsilon > 0$ a number $\delta > 0$ can be given with which for any $x \in \mathcal{D}_{f}, |x-a| < \delta$ we have $|f(x)-f(a)| < \varepsilon$. \\
Notation: $f \in \mathcal{C}\{a\}$; if $\forall x \in \mathcal{D}_{f} : f \in \mathcal{C}\{x\}$, then $f \in \mathcal{C}$.

\section{Relationship between limit and continuity}

Let $f \in \mathbb{K}_{1} \to \mathbb{K}_{2}, a \in \mathcal{D}_{f} \cap \mathcal{D}'_{f}$. Then $f \in \mathcal{C}\{a\} \iff \lim\limits_{a}{f} = f(a)$.

\section{Concepts}

\subsection{Compact set}

$A \subset \mathbb{K}$ is compact if any sequence $(x_{n}): \mathbb{N} \to A$ has a subsequence $(x_{\nu_{n}})$ that is convergent and $\lim{(x_{\nu_{n}})} \in A$. \\
Then $A$ is bounded and closed.

\subsection{Convergent}

A number sequence $x = (x_{n}) : \mathbb{N} \to \mathbb{K}$ is called convergent if $\exists \alpha \in \mathbb{K}, \forall \varepsilon > 0, \exists N \in \mathbb{N}, \forall n \in \mathbb{N}, n > N : |x_{n}-\alpha| < \varepsilon$. \\
$\alpha$ is the limit of the sequence $x$.

\subsection{Bounded}

For a sequence: $x_{n}$ bounded $\Rightarrow \exists \nu : x \circ \nu$ convergent \\
For a set: $\emptyset \neq A \subset \mathbb{K}$ is bounded if $\exists q \in \mathbb{R} : |x| \leq q, (x \in A)$

\subsection{Closed set}

Its complement is an open set.

\subsection{Open set}

$A$ is an open set $\iff \forall a \in A, \exists K(a): K(a) \subset A$, that is, every point of the set $A$ is an interior point.

\section{Heine theorem {\large (optional)}}

If the function $f \in \mathbb{K}_{1} \to \mathbb{K}_{2}$ is continuous and $\mathcal{D}_{f}$ is compact, then $f$ is uniformly continuous.

\subsection{Uniformly continuous}

The function $f \in \mathbb{K}_{1} \to \mathbb{K}_{2}$ is uniformly continuous if $\forall \varepsilon > 0, \exists \delta > 0 : |f(x)-f(t)| < \varepsilon, (x,t \in \mathcal{D}_{f}, |x-t| < \delta)$.

\section{Weierstrass theorem}

Suppose that the function $f \in \mathbb{K} \to \mathbb{R}$ is continuous and $\mathcal{D}_{f}$ is compact. Then the range $\mathcal{R}_{f}$ has a greatest and a smallest element $(\exists \max{\mathcal{R}_{f}}, \exists \min{\mathcal{R}_{f}})$.

\section{Bolzano theorem}

Suppose that for some $-\infty < a < b < +\infty$ the function $f:[a,b] \to \mathbb{R}$ is continuous, and $f(a) \cdot f(b) < 0$ (of opposite sign). \\
Then there exists $\xi \in (a,b)$ such that $f(\xi) = 0$.

\section{The concept of a power series}

Let the center $a \in \mathbb{K}$ and the coefficient sequence $(a_{n}): \mathbb{N} \to \mathbb{K}$ be given, and using these consider the following functions: $f_{n}(t) := a_{n}(t-a)^{n}, (t \in \mathcal{D} := \mathbb{K}, n \in \mathbb{N})$. Then the function series $\sum{(f_{n})}$ is called a power series.

\subsection{Sequence}

The function $f$ is called a sequence if $\mathcal{D}_{f} = \mathbb{N}$.

\subsection{Function sequence, function series}

A sequence $(f_{n})$ is a function sequence if for $\forall n \in \mathbb{N}$ each $f_{n}$ is a function, and for some set $\mathcal{D} \neq \emptyset$ we have $\mathcal{D}_{f_{n}} = \mathcal{D}, (n \in \mathbb{N})$. \\
If for the function sequence $(f_{n})$ in question $\mathcal{R}_{f_{n}} \subset \mathbb{K}, (n \in \mathbb{N})$ also holds, then the function series $\sum{(f_{n})}$ determined by the function sequence $(f_{n})$ is the following function sequence: $\sum{(f_{n})} := (\sum\limits_{k=0}^{n}{f_{k}})$.

\section{Cauchy--Hadamard theorem}

Suppose that for the sequence $(a_{n}) : \mathbb{N} \to \mathbb{K}$ the limit $\lim{(\sqrt[n]{|a_{n}|})}$ exists, and let
$$ r := \left\{\begin{array} {lr}
        +\infty                              & \text{if } \lim{(\sqrt[n]{|a_{n}|})} = 0 \\
        \dfrac{1}{\lim{(\sqrt[n]{|a_{n}|})}} & \text{if } \lim{(\sqrt[n]{|a_{n}|})} > 0
    \end{array}\right.$$
We call $r$ the radius of convergence.
Then for any $a \in \mathbb{K}$ we can say the following about the power series $\sum{(a_{n}(t-a)^{n})}$:
\begin{enumerate}
    \item If $r > 0$, then at every point $x \in \mathbb{K}, |x-a| < r$ the power series $\sum{(a_{n}(t-a)^{n})}$ is absolutely convergent at $x$.
    \item If $r < +\infty$, then for any $x \in \mathbb{K}, |x-a| > r$ the power series $\sum{(a_{n}(t-a)^{n})}$ is divergent at $x$.
\end{enumerate}

\subsection{Absolutely convergent}

The infinite series $\sum\limits_{n=1}^{\infty}{a_{n}}$ is called absolutely convergent if the series $\sum\limits_{n=1}^{\infty}{|a_{n}|}$ is convergent.

\subsection{Divergent}

If $\lim\limits_{n \to \infty}{a_{n}}$ does not exist, or is not finite, then the infinite series $\sum{a_{n}}$ is divergent.

\subsection{Consequences of the Cauchy--Hadamard theorem}

\begin{enumerate}
    \item If $r = +\infty$, then $\mathcal{D}_{0} = \mathbb{K}$, and for $\forall x \in \mathbb{K}$ the power series is absolutely convergent.
    \item If $r = 0$, then $\mathcal{D}_{0} = \{a\}$ and $\sum\limits_{n=0}^{\infty}{a_{n}(a-a)^{n}} = a_{0}$.
    \item If $0 < r < +\infty$, then $K_{r}(a) \subset \mathcal{D}_{0} \subset G_{r}(a) = \{x \in \mathbb{K} : |x-a| \leq r\}$ and at every point $x$ of the neighborhood $K_{r}(a)$ the power series is absolutely convergent at $x$.
    \item If $r > 0$, then for any number $0 \leq \rho < r$ choose the element $x \in K_{r}(a)$ such that $|x-a| = \rho$. Then $\sum\limits_{n=0}^{\infty}{|a_{n}(x-a)^{n}|} = \sum\limits_{n=0}^{\infty}{|a_{n}||x-a|^{n})} = \sum\limits_{n=0}^{\infty}{|a_{n}|\rho^{n})} < +\infty$
\end{enumerate}

\section{Analytic functions}

Suppose that the radius of convergence $r$ (see C--H theorem) of the power series $\sum{(a_{n}(t-a)^{n})}$ is not zero. Then we can define the following function: $f(x) := \sum\limits_{n=0}^{\infty}{a_{n}(x-a)^{n}}, (x \in K_{r}(a))$, which is nothing other than the restriction of the sum function $F(x) := \sum\limits_{n=0}^{\infty}{a_{n}(x-a)^{n}}, (x \in \mathcal{D}_{0})$ of the function series $\sum{(a_{n}(t-a)^{n})}$ to the neighborhood $K_{r}(a)$: $f = F|_{K_{r}(a)}$, ($f$ is the sum function of the power series). A function $f$ of such a structure is called analytic. \\
Remark: $\mathcal{D}_{0}$ denotes the domain of convergence of the power series $\sum{(a_{n}(t-a)^{n})}$.

\subsection{Important analytic functions}

Coefficients $a_{n}$ for which $\lim{(\sqrt[n]{|a_{n}|})} = 0$, hence $r = +\infty$ for the power series $\sum{(a_{n}t^{n})}$. \vspace{0.3cm} \\
These $a_{n}$ are: $\dfrac{1}{n!}, \dfrac{(-1)^n}{(2n+1)!}, \dfrac{(-1)^n}{(2n)!}, \dfrac{1}{(2n+1)!}, \dfrac{1}{(2n)!}$.
\begin{itemize}
    \item $\exp{x} := \exp{(x)} = e^{x} = \sum\limits_{n=0}^{\infty}{\dfrac{x^n}{n!}, (x \in \mathbb{K})}$
    \item $\sin{x} := \sin{(x)} = \sum\limits_{n=0}^{\infty}{(-1)^n\dfrac{x^{2n+1}}{(2n+1)!}, (x \in \mathbb{K})}$
    \item $\cos{x} := \cos{(x)} = \sum\limits_{n=0}^{\infty}{(-1)^n\dfrac{x^{2n}}{(2n)!}, (x \in \mathbb{K})}$
    \item $\sinh{x} := \sinh{(x)} = \sum\limits_{n=0}^{\infty}{\dfrac{x^{2n+1}}{(2n+1)!}, (x \in \mathbb{K})}$
    \item $\cosh{x} := \cosh{(x)} = \sum\limits_{n=0}^{\infty}{\dfrac{x^{2n}}{(2n)!}, (x \in \mathbb{K})}$
\end{itemize}

\end{document}
